\section{Képviseleti Demokrácia}

Ez a fejezet a képviseleti demokráciát és hiányosságait tárgyalja. Először a szükséges
fogalmakat definiáljuk, majd a rendszer alapvető hibát részletezzük erőssen
támaszkodva korábbi intellektuális kritikákra, ezután megvizsgáljuk a rendszer működését
egy ideális környezetben és belátjuk, hogy ebben az esetben sem produkálja az elvárt
eredményt ez a megközelítés, végezetül pedig azt vizsgáljuk, hogyan --- és hogyan nem ---
működik a képviseleti demokrácia a valóságban.

A Britanicca a következőképpen definiálja a képviseleti demokráciát\cite{rep_dem_def} és
a továbbiakban mi is ezt fogjuk használni (fordítás):

\begin{displayquote}
egy politikai rendszer, amelyben egy ország vagy más politikai egység polgárai képviselőkre
szavaznak, akik a nevükben végzik a törvényhozást és irányítják az adott egységet. A
megválasztott képviselők viszont elszámoltathatók a választópolgárok által.
A demokrácia egyik formájaként a képviseleti demokrácia a közvetlen demokráciával áll szemben,
amelyben minden állampolgár közvetlenül szavaz a törvényekről és egyéb kérdésekről. A legtöbb
modern ország képviseleti demokrácia, és ennek megfelelően számos kihívással néz szembe.
\end{displayquote}

Általánoságban a képviseleti demokrácia a következőképpen működik: egy adott választókerület
polgárai kiválasztják azt a személyt, aki szerintük a legjobban képviselni tudja az érdekeiket.
A mandátum megszerzése után a képviselő --- limitált korlátokkal --- a belátása szerint szavaz
a kormányban és alkot saját előterjesztéseket. Más szóval, a demokrácia ezen formájának a
feltételezése, hogy a képviselő rendelkezik az emberek beleegyezésével a következő néhány évre.
Ha kritikusan elmezzük ezt a kijelentést akkor észrevehetjük, hogy ez igen szokatlan. Ami azt
illeti életünk egyik másik aspektusában sem fogadnánk el ezeket a feltételeket. Gondoljunk bele
mennire máshogy működne a házasság intézménye ha a mennyasszony életéről 5 évig a párja döntene,
legyen utóbbi akármilyen aggresszív vagy hozzon bármilyen rossz döntést. A képviseleti demokrácia
támogatói valószínűleg amelett érvelnének, hogy a fékek és ellensúlyok megakadályozzák az effajta
túlkapásokat. Azonban ez csak délibáb. Az előző példához visszatérve, a férj nem csak a
redőrkapitányt és a bírókat nevezte ki, de jó kapcsolatot ápol a helyi TV csatorna vezetőjével is
így --- talán nem meglepő módon --- a városban mindenki tudja, hogy a feleség sportolás közben 
törte el a kezét és amúgy sem elég hálás az urának. A progresszív olvasók valószínűleg tudják,
hogy az előbb említettek nagyon is hasonlítanak arra ahogy a házasság működött nem is olyan régen
és hogy a nők keserves küzdelem árán változtattak ezen a helyzeten. A változás pozitív hatásait ma
már tudomány is alá támasztja \cite{divorce}. A gyengébbik nem megtanulta, hogy nem mondhat le az
autonomiáról és nem adhat feltétlen beleegyezést párjának az élete felett akármilyen ígéretes is a
kérő. Egy lecke ami hasznos lehet az emberiség egészének, az állam működésének újra gondolásához.  

Ezen a ponton fontos megjegyezni, a félreértések elkerülése végett, hogy nem amellett érvelünk,
hogy \textit{fékek és ellensúlyok nem működnek}, hanem sokkal inkább amelett, hogy \textit{fékek
és ellensúlyok egy olyan homogén csoport kezében vannak, akiknek nem érdekük, hogy ezen eszközök
hatékonyak legyenek}. Néhány kivétellel a képviselők hasonló gazdasági és szociális háttérrel
rendelkeznek \cite{reps_are_homogen_1, reps_are_homogen_2}, és politikai iranyultságtól függetlenül
ugyanazon pénzügyi lehetőségek állnak rendelkezésükre a mandátum szerzés után. Erre egy jó példa az
Egyesült Államok esete, ahol ugyan léteznek törvények amelyek arra hivatottak, hogy megakadályozzák
a tőzsdei nyerészkedést belső információk alapján, azonban ezek nem kerülnek betartásra 
\cite{insider_trading}. Természetesen, se a Republikánus se a Demokrata pártnak nem érdeke ezen
szabályok betartatása, hiszen akkor maguk sem profitálhatnának a szabályok be nem tartásából.
Ilyen és ehhez hasonló példák mutatják, hogy az önszabályozás nem biztosít garanciát a törvények
betartására.

Az önszabályozás eredménytelensége azonban csak egy a megannyi hibából, amit az értelmiség
megfogalmazott az utóbbi néhány évtizedben. A képviseleti demokrácia kritikája szerint a
rendszer hajlamos arra, hogy a politikai döntéshozatal egy szűk elit kezében összpontosuljon,
miközben a választók szerepe nagyrészt a versengő elitcsoportok közötti választásra korlátozódik
\cite{rep_dem_critic_1, rep_dem_critic_3}. Ideális esetben a demokratikus rendszerek jól tükröznék
a nép preferenciáit, a gyakorlatban a politikai folyamatokat gyakran intézményi és strukturális
korlátok formálják amelyek eltávolíthatják a döntéshozatalt az állampolgároktól \cite{rep_dem_critic_2}.
A modern társadalmakban továbbá megfigyelhető, hogy a politikai részvétel gyengül, a pártok és
intézmények elszakadnak a társadalomtól, és a politikai verseny egyre inkább professzionális
elitjáték lesz \cite{rep_dem_critic_4, rep_dem_critic_5}. Emellett empirikus kutatások arra utalnak,
hogy a választók döntései gyakran nem stabil preferenciákon vagy mély politikai ismereteken alapulnak,
hanem rövid távú hatások és torzítások befolyásolják őket, ami tovább gyengíti a képviseleti rendszer
demokratikus minőségét \cite{rep_dem_critic_6}.

Talán a képviseleti demokrácia védelmezői amellett érvelnének, hogy nem a rendszerrel magával van
a probléma, hanem sokkal inkább a szélesebb szocio-gazdasági környezettel, ami korrupcióhoz vezet
a politikai résztvevők körében. A következő fejezetben bemutatjuk, hogy a demokrácia ezen formája
ideális esetben sem működik, mivel nem képes elég részletességében reprezentálni a népakaratot.

\subsection{Ideális Világ, Nem Idális Rendszer}

Ebben a fejezetben azt mutatjuk be, hogy egy ideális rendszerben sem működik a képviseleti
demokrácia, mivel képviselők nem képesek hatékonyan reprezentálni az összes kerületükhöz
tartozó választót minden egyes kérdésben. Ennek oka, hogy egy választópolgárok preferenciái
komplexek és gyakran nem kategorizálhatóak könnyen jól definiálható kategoriákba. Mivel a
demokrácia ezen formájának nincs eszköze, amivel könnyen, gyorsan, és költséghatékonyan a polgárok
kifejezhetik véleményüket bármely előterjesztésről, ennek eredménye jogszabályi stagnálás és
törvények amivel a többség nem ért egyet.

Egy másik próbléma a jelenlegi rendszerrel, hogy ugyan elméletben bárki elindulhat a választásokon
a gyakorlatban a legtöbb képviselő elit osztályból kerül ki, gyakran ugyan azon családokból
generációkon át \cite{reps_are_homogen_1, reps_are_homogen_2, reps_are_homogen_3}. A képviseleti
demokrácia alapfeltevése az, hogy ha elég okos, képzett és tapasztalt embereket választunk meg
döntéshozónak, akkor jó döntések születnek. A probléma ezzel --- feltételezve, hogy ezen egyének
nem ostobák vagy rosszindulatúak ---, hogy túl hasonlók egymáshoz. A társadalomtudomány ezt a
jelenséget homofíliának nevezi: az emberek hajlamosak azokkal együtt gondolkodni és dönteni, akik
hozzájuk hasonló háttérrel, képzettséggel és világlátással rendelkeznek. Ez pedig nem független
hibákhoz, hanem korrelált tévedésekhez vezet. A döntéstudomány egyik legismertebb példája ezt jól
illusztrálja. Francis Galton 1906-ban \cite{wisdom_of_crowds, wisdom_of_crowds_wiki} egy vásáron
megkérdezte az embereket, mennyit nyom egy ökör. Az egyéni becslések nagy része pontatlan volt,
de az átlaguk meglepően közel esett a valós értékhez. Ennek oka, hogy a résztvevők különböző
háttérrel, eltérő intuíciókkal rendelkeztek, így a hibáik kioltották egymást. A képviseleti
demokráciában azonban nem ez történik. A politikai elit jellemzően ugyanazokból a társadalmi
rétegekből kerül ki, hasonló iskolákba jártak, ugyanazokat az intézményeket tisztelik. Így a
hibáik nem kioltják, hanem megerősítik egymást. Ezt támasztja alá Lu Hong és Scott Page
\cite{diverse_groups_are_better} kutatása is, amely kimutatta, hogy homogén, 
\textit{legjobb szakértőkből} álló csoportok rosszabbul teljesítenek komplex problémák megoldásában,
mint sokszínű csoportok, ahol eltérő nézőpontok ütköznek.

A képviseleti rendszer másik alapvető hibáját a második világháború egyik legismertebb statisztikai
története világítja meg. Az amerikai hadsereg a visszatérő bombázók sérüléseit vizsgálva akarta
eldönteni, hova kell páncélt erősíteni. Abraham Wald \cite{aircraft} mutatott rá, hogy pont azokon
a helyeken kell megerősíteni a gépeket, ahol nincs sérülés --- mert azok a gépek, amelyeket ott
találtak el, nem tértek vissza. A döntéshozók csak a túlélőket látták, és majdnem végzetes következtetést
vontak le. Ugyanez történik a képviseleti demokráciában is. A parlamentekben, minisztériumokban és
szakértői testületekben csak a túlélők hangja hallatszik: az eliteké, az intézményeken belül mozgóké. A
rendszerből kiszorulók, a marginalizált csoportok, a mindennapi tapasztalatból fakadó tudás nem jut
be a döntéshozatalba. Így a konszenzus nem az igazság jele lesz, hanem a konformitásé. 

Ezért nem működik a képviseleti demokrácia: nem azért, mert az emberek eredendően rosszak, hanem mert
egy olyan struktúrát hoz létre, amely rendszerszinten vak a saját hibáira. Ennélfogva a rendszer 
alapvető hibája elkerülhetetlen, még egy olyan világban is ahol nincsenek egyenlőtlenségek és minden
képviselő a moralitás megtestesítője. A tudomány világosan mutatja, hogy komplex társadalmi problémákat
nem lehet szűk, homogén elit kezében hagyni. Valódi demokráciához nem jobb képviselőkre, hanem szélesebb,
közvetlenebb részvételre és sokszínű döntéshozatalra van szükség.

Ez a fejezet rendkívül kritikusan fogalmazott a képviseleti demokráciáról, de fontos megjegyezni, hogy
ez a rendszer egy olyan világban született, ahol távkommunikáció még nem létezett. Egy ilyen világban
a képviseleti demokrácia egy megfelelő kompromisszum volt. Azonban manapság aggregálni és feldolgozni
milliónyi véleményt nem csak lehetséges, hanem a mindennappok része, legyen ez közösségi média vagy
globális piaci analízis. A technológia létezik, ideje használni.

\subsection{Nem Ideális Világ, Nem Idális Rendszer}

Eddig nem érintettük a korrupció kérdését, ennek két oka volt. Az első, hogy ez a kérdés alaposan
tanulmányozott a szakértők és a média által \cite{corruption} ennél fogva általánoságban véve közismert.
A második pedig, hogy értekezések a korrupcióról hajlamosak egyénekre hárítani a felelősséget, a rendszer
helyett ami valódi okozója ezeknek a problémáknak.

Ezek ellenére, a korrupció egy fontos kérdes minden országban, ami negatívan befolyásolja valamennyiünk
életet \cite{corruption_and_inequality}. Azonban fontos belátni, hogy ez a probléma elkerülhetetlen
a jelenlegi rendszerben. A képviselők túl sok hatalommal bírnak a mandátumok megszerzése után és az
elszámoltatás túl későn --- gyakran csak a következő választáskor --- vagy sosem történik. Azonban
mindenki a bőrén érzi a korrpuciót. Legyen ez az ingatlan árak mesterségesen magasan tartása, csakhogy
néhány, tucatnyi ingatlannal rendelkező politikus még gazdagabbá váljon. Vagy megfigyelő szoftverek
telepítése mindenki telefonjára a gyermek védelem nevében, csak mert a szoftvert biztosító cég hozzájárult
bozonyos EU-s képviseleők kampányához. Akármilyen formában is a korrupció előbb-utóbb mindannyiunk életét
megkeseríti.  

A képviseleti demokrácia hívei valószínűleg amelett érvelnének, hogy ezt a problémát megfelelő
jogszabályokkal orvosolni lehet. Azonban akármilyen átgondolt és mindent lefedő törvényeket hozunk,
ezen jogszabályok betartása akkor is a hatalmon lévőkre és a kinevezettjeikre hárul. Ez a rendszer
nem ajánl garanciát a korrupció felszámolására, csak reményt, hogy a következő képviselők a helyes
döntést fogják hozni, egy olyan környezetben ahol túl könnyű a rosszul cselekedni. A remény nem
lehet alapkő, amire egy egész társadalmat építünk. Amire szükség van az az, hogy a választók valódi
kontrollt tudjanak kifejteni a törvénykezésre és ha szükséges az általuk megfelelő irányba tereljék
azt. Nem négyévente, ha van kihívó akinek ideje és vagyona engedi, hogy elinduljon a választáson,
hanem minden nap.

Ezzel el is érkeztünk a cikk fő témájához, a valódi alternatívákhoz. A következő fejezetben amellett
érvelünk majd, hogy a direkt demokrácia képes medoldást nyújtani megannyi problémára amit a
képviseleti demokrácia okoz.

